\documentclass[letterpaper,12pt]{article}
\usepackage{graphicx}
\usepackage{pdflscape}
\usepackage{amsmath}
\usepackage{amsthm}
\usepackage{lscape}
\usepackage{enumerate}
\usepackage{float}
\usepackage{grffile}
\usepackage{relsize}
\usepackage{tablefootnote}
\usepackage{footnote}
\usepackage{sectsty}
\usepackage{xcolor}
\usepackage{caption}
\usepackage{color}
\usepackage{natbib}
\usepackage{multirow}
\usepackage{bbm}
\usepackage[T1]{fontenc}
\usepackage[multiple]{footmisc}
\usepackage{appendix}
\usepackage{times}
\usepackage{relsize}
\usepackage{caption}
\usepackage{newtxtext,newtxmath}
\usepackage{enumitem}
\usepackage{booktabs}
\usepackage{tikz}
\usetikzlibrary{arrows.meta,calc,intersections}

\definecolor{sol}{RGB}{0,80,160}
\definecolor{emp}{RGB}{180,0,0}

\title{ECON311: Assignment II --- Solutions}
\author{}
\date{Fall 2026}

\newlength{\defbaselineskip}
\setlength{\defbaselineskip}{\baselineskip}
\newcommand{\setlinespacing}[1]{\setlength{\baselineskip}{#1 \defbaselineskip}}
\newcommand{\doublespacing}{\setlength{\baselineskip}{1.3 \defbaselineskip}}
\newcommand{\singlespacing}{\setlength{\baselineskip}{1\defbaselineskip}}
\renewcommand{\thesubsection}{\arabic{subsection}}

\textwidth=6.5in
\textheight=9in
\oddsidemargin=0.10in
\evensidemargin=0.10in
\topmargin=-0.5in

\begin{document}
\pagenumbering{arabic}
\maketitle
\vspace{-.4in}


%----------------------------------------------------------------------
\section*{Question I: Long-Run Growth [12 points]}
%----------------------------------------------------------------------

\noindent Model: $Y_t = A_t L_{Y_t}$,\quad $A_{t+1} = A_t(1 + \overline{z}\,\overline{\ell}\,\overline{L})$,\quad $L_{A_t} = \overline{\ell}\,\overline{L}$,\quad $L_{Y_t} = (1-\overline{\ell})\,\overline{L}$.

\subsection*{(a) Growth Rate and the R\&D Policy Shock [6 points]}

\noindent\textbf{Deriving $g_y$:}

Output per person: $y_t = Y_t/\overline{L} = A_t(1-\overline{\ell})$.

Since $(1-\overline{\ell})$ is constant, $g_y = g_A$. From the law of motion:
\[
\frac{A_{t+1}}{A_t} = 1 + \overline{z}\,\overline{\ell}\,\overline{L} \qquad \Longrightarrow \qquad \textcolor{sol}{g_y = g_A = \overline{z}\,\overline{\ell}\,\overline{L}}
\]

\noindent\textbf{i. Effect on $g_y$ when $\overline{\ell}$ increases to $\overline{\ell}' > \overline{\ell}$:}

\textcolor{sol}{$g_y' = \overline{z}\,\overline{\ell}'\,\overline{L} > \overline{z}\,\overline{\ell}\,\overline{L} = g_y$.}

The growth rate rises \textbf{immediately and permanently}. This is a \emph{growth rate effect}.

\noindent\textbf{ii. Effect on the level of $y_t$:}

At $t_0$:
\[
y_{t_0} = A_{t_0}(1-\overline{\ell}') < A_{t_0}(1-\overline{\ell})
\]
Since $\overline{\ell}' > \overline{\ell}$, the level of $y$ \textbf{falls at $t_0$}: fewer workers are now producing goods (more are in R\&D). This is the short-run cost.

In the long run, however, $A_t$ grows at the faster rate $\overline{z}\,\overline{\ell}'\,\overline{L}$, so $y_t$ eventually overtakes the old trajectory and rises above it permanently. There is therefore both a (negative) level effect at $t_0$ and a positive growth rate effect afterward.

\noindent\textbf{iii. Graph on ratio scale ($\ln y_t$ vs.\ time):}

The old balanced growth path (BGP) is a straight line with slope $g_y$. At $t_0$, $\ln y$ \textbf{drops} (discrete downward jump of size $\ln(1-\overline{\ell}') - \ln(1-\overline{\ell}) < 0$). From $t_0$ onward, $\ln y$ rises along a new line with the steeper slope $g_y' > g_y$. The new trajectory starts below the old one but crosses it at some finite time and lies above it thereafter. The graph should clearly label: (1) the level drop at $t_0$, (2) the steeper post-$t_0$ slope, and (3) the eventual upward crossing.

\subsection*{(b) Two Countries with Different $\bar{z}$ [6 points]}

Since $g_y^1 = \overline{z}^1\overline{\ell}\,\overline{L} > \overline{z}^2\overline{\ell}\,\overline{L} = g_y^2$ and both start from the same $\ln y$ in 1980:

\noindent\textbf{i. Graph:} Both lines start at the same point in 1980. Country~1's line has a strictly steeper slope. From 1980 onward, country~1 has higher output per person for all $t > 1980$, with the gap growing over time.

\noindent\textbf{ii. Long-run gap:}

\textcolor{sol}{The gap in log output per person grows without bound:}
\[
\ln y_t^1 - \ln y_t^2 = (g_y^1 - g_y^2)\,t \to \infty \text{ as } t \to \infty
\]

The gap \textbf{widens} over time and never converges.

\textbf{Contrast with Solow model:} The Solow model predicts \emph{conditional convergence} — countries with the same parameters converge to the same steady state. If two Solow economies have the same $\overline{s}$, $\overline{d}$, and $\overline{A}$, income gaps narrow regardless of initial conditions. Here, different $\overline{z}$ leads to permanently different growth rates, so income disparities compound indefinitely. Countries that differ only in their capacity to generate ideas face ever-widening inequality. This is a striking and empirically relevant implication: innovation policy has permanent, compounding consequences for relative living standards.


%----------------------------------------------------------------------
\section*{Question II: The IS-MP Model and the Lower Bound on Interest Rates [10 points]}
%----------------------------------------------------------------------

\subsection*{(a) The ZLB Constraint [5 points]}

\noindent\textbf{Setting:} IS-MP model with IS: $\widetilde{Y} = \bar{a} - b(R-\bar{r})$ and MP rule: $R = \bar{r} + \bar{m}(\pi-\bar{\pi})$.

Large negative demand shock: $\bar{a}$ falls sharply to $\bar{a}' \ll 0$, shifting IS far to the left.

To restore $\widetilde{Y}=0$, the CB must reduce $R$ to:
\[
R^* = \bar{r} + \frac{\bar{a}'}{b} < \bar{r} \quad (\text{since } \bar{a}' < 0)
\]

If $R^* < R_{\min} = -\pi$, the CB cannot reach the target. At $R = R_{\min}$:
\[
\widetilde{Y}_{ZLB} = \bar{a}' - b(R_{\min} - \bar{r}) < 0
\]

\textbf{IS-MP diagram:} The original IS and MP cross at $(\widetilde{Y}=0,\, R=\bar{r})$. After the shock, IS' lies far to the left. The CB lowers $R$ along the vertical axis until hitting $R_{\min}$ — shown as a horizontal floor on the diagram. The intersection of IS' with the $R_{\min}$ floor is at $\widetilde{Y}_{ZLB} < 0$. The CB \emph{cannot} lower $R$ further to close the remaining gap.

\textbf{Why does the ZLB bind?} Because nominal interest rates are bounded below at zero (lenders will not pay to lend when they can hold cash), real rates face the floor $R \geq -\pi$. When a large recession would require deeply negative real rates to stimulate demand, monetary policy loses its conventional lever.

\textbf{Policy alternatives:} Fiscal policy (government spending, tax cuts), forward guidance (committing to future low rates), quantitative easing (CB purchases long-dated assets to lower long-term rates), or unconventional monetary tools.

\subsection*{(b) Fiscal Policy Effectiveness at the ZLB [5 points]}

At the ZLB, $R = R_{\min}$ is fixed; the CB's rule is non-binding. A fiscal expansion $\Delta\overline{G} > 0$ shifts IS' rightward to IS''.

\textbf{IS-MP diagram:} IS'' lies to the right of IS', intersecting the fixed $R_{\min}$ floor at a higher $\widetilde{Y}$, closer to zero.

\textbf{Why is the fiscal multiplier larger at the ZLB?}

In \textbf{normal times}, the sequence is:
\[
\Delta\overline{G} \uparrow \;\to\; \widetilde{Y} \uparrow \;\to\; \pi \uparrow \;\to\; R \uparrow \text{ (CB rule)} \;\to\; I \downarrow \;\to\; \widetilde{Y} \downarrow \text{ (partial crowding out)}
\]
The CB's automatic tightening partially offsets the fiscal stimulus. The effective multiplier is \emph{smaller} than the pure Keynesian multiplier.

At the \textbf{ZLB}, the CB cannot raise $R$ (already at the floor). Therefore:
\[
\Delta\overline{G} \uparrow \;\to\; \widetilde{Y} \uparrow \;\to\; \pi \uparrow \;\to\; R \text{ is unchanged} \;\to\; \text{no crowding out}
\]
The full Keynesian multiplier $1/(1-b_c)$ applies, undiluted by any monetary offset. Formally, from the IS equilibrium at fixed $R$:
\[
\widetilde{Y} = \frac{\Delta\overline{G}}{1-b_c} \times \frac{1}{\overline{Y}}
\]
versus the normal-times IS-MP combined multiplier, which is strictly smaller. The ZLB therefore makes fiscal policy unusually powerful — a key lesson from the 2008 financial crisis and the 2020 pandemic recession.


%----------------------------------------------------------------------
\section*{Question III: The Investment Accelerator, IS-MP and Automatic Stabilisers [35 points]}
%----------------------------------------------------------------------

\noindent Model: $C = a_c + b_c(Y-T)$, $T = b_T Y$, $I = a_i + b_a Y - b_i R$, $G = \overline{G}$, $NX = \overline{X}-\overline{M}$.

\subsection*{(a) Goods-Market Equilibrium and Multiplier [4 points]}

Aggregate expenditure:
\begin{align*}
E &= a_c + b_c(1-b_T)Y + a_i + b_a Y - b_i R + \overline{G} + \overline{X} - \overline{M} \\
  &= \underbrace{\bigl[a_c + a_i - b_i R + \overline{G} + \overline{X} - \overline{M}\bigr]}_{\text{autonomous}} + \underbrace{\bigl[b_c(1-b_T) + b_a\bigr]}_{\text{induced}}Y
\end{align*}

Setting $Y = E$:
\[
Y\bigl[1 - b_c(1-b_T) - b_a\bigr] = a_c + a_i - b_i R + \overline{G} + \overline{X} - \overline{M}
\]

\textcolor{sol}{
\[
Y = \frac{a_c + a_i - b_i R + \overline{G} + \overline{X} - \overline{M}}{1 - b_c(1-b_T) - b_a}, \qquad \text{Multiplier} = \frac{1}{1-b_c(1-b_T)-b_a}
\]
}

\textbf{Stability condition:} $b_c(1-b_T) + b_a < 1$ (so denominator is positive and finite).

\textbf{Comparison with $b_a = 0$:} Without the accelerator, multiplier $= 1/[1-b_c(1-b_T)]$. With $b_a > 0$, the denominator is smaller $\Rightarrow$ multiplier is \textbf{larger}. The accelerator amplifies each round of spending: higher $Y$ triggers additional investment ($b_a Y$), which raises $Y$ further.

\subsection*{(b) Keynesian Cross with an Investment Shock [6 points]}

An increase of $\Delta a_i$ shifts the expenditure line \textbf{upward by $\Delta a_i$} (parallel shift).

\textbf{With accelerator ($b_a > 0$):} The slope of the expenditure line is $b_c(1-b_T) + b_a$ — steeper than without the accelerator. A given upward shift $\Delta a_i$ produces a larger $\Delta Y$:
\[
\Delta Y = \frac{\Delta a_i}{1-b_c(1-b_T)-b_a} > \frac{\Delta a_i}{1-b_c(1-b_T)} = \Delta Y\big|_{b_a=0}
\]

\textbf{Economic mechanism:} Investment rises by $\Delta a_i$ (impact). This raises $Y$ (round 1). Higher $Y$ induces further investment through the accelerator $b_a \Delta Y$ (round 2). This raises $Y$ again, inducing yet more investment (round 3). Each round is smaller, and the process converges if the stability condition holds.

\textbf{Keynesian Cross diagram:} Expenditure line has slope $b_c(1-b_T)+b_a$ (steep) for the accelerator case and slope $b_c(1-b_T)$ (flatter) for the no-accelerator case. Both shift up by $\Delta a_i$ from the same initial equilibrium. The new $Y$ is higher under the accelerator, as the intersection with the 45° line is further to the right.

\subsection*{(c) The IS Curve [5 points]}

At potential output $\overline{Y}$ (with $R = \overline{r}$):
\[
\overline{Y} = \frac{a_c + a_i - b_i\overline{r} + \overline{G} + \overline{X}-\overline{M}}{1-b_c(1-b_T)-b_a}
\]

For general $R$:
\[
\widetilde{Y} = Y - \overline{Y} = \frac{-b_i(R-\overline{r})}{1-b_c(1-b_T)-b_a} \equiv -\tilde{b}(R-\overline{r})
\]

\textcolor{sol}{
\[
\text{IS curve:} \quad \widetilde{Y} = -\tilde{b}(R-\overline{r}), \qquad \tilde{b} \equiv \frac{b_i}{1-b_c(1-b_T)-b_a}
\]
}

\textbf{Effect of higher $b_a$:} Since $1-b_c(1-b_T)-b_a < 1-b_c(1-b_T)$, we have $\tilde{b} > b_i/[1-b_c(1-b_T)]$. The IS curve is \textbf{flatter} (a given $\Delta R$ produces a larger $\Delta\widetilde{Y}$). Economic interpretation: higher $R$ reduces $I$ directly; through the accelerator, this fall in $I$ reduces $Y$, which further reduces $I$, amplifying the effect of the interest rate change.

\subsection*{(d) IS-MP-PC Dynamics after an Investment Boom [9 points]}

Let $t$ and $t+1$ be the two periods of the boom ($a_i$ elevated), and $t+2$ onward the return to normal ($a_i$ at original level). Start at LR equilibrium: $\widetilde{Y}=0$, $\pi=0.03$, $R=\overline{r}$.

\noindent\textbf{Period $t$:} IS shifts right. $\widetilde{Y}_t > 0$.
\[
\pi_t = \pi_{t-1} + \overline{v}\,\widetilde{Y}_t = 0.03 + \overline{v}\,\widetilde{Y}_t > 0.03
\]
CB raises $R_t$ in response to rising $\pi_t$ (MP rule).

\noindent\textbf{Period $t+1$:} IS still shifted right (boom continues). Higher $\pi_t$ has induced the CB to raise $R$ above $\overline{r}$, partially restraining the boom. $\widetilde{Y}_{t+1} > 0$ but smaller than $\widetilde{Y}_t$ (since $R$ is higher). Inflation rises further: $\pi_{t+1} > \pi_t > 0.03$.

\noindent\textbf{Period $t+2$:} Boom ends; IS returns to original position. Inflation is now elevated above target. The CB, following its rule, keeps $R > \overline{r}$ because $\pi > \overline{\pi}$. This pushes $\widetilde{Y}_{t+2} < 0$ (mild recession).
\[
\pi_{t+2} = \pi_{t+1} + \overline{v}\,\widetilde{Y}_{t+2} < \pi_{t+1}
\]

\noindent\textbf{Subsequent periods:} The negative output gap gradually pulls inflation back toward 0.03\%. As $\pi$ falls toward $\overline{\pi}$, $R$ falls toward $\overline{r}$, $\widetilde{Y}$ returns toward 0. Economy converges back to LR equilibrium.

\noindent\textbf{Summary of time paths:}
\begin{itemize}
  \item $\widetilde{Y}$: rises to positive during boom, falls below zero after boom ends, converges back to 0.
  \item $\pi$: rises above 3\% during boom, continues rising as $\widetilde{Y}>0$; slowly falls back to 3\% as $\widetilde{Y}<0$ during the post-boom correction.
\end{itemize}

\noindent\textbf{IS-MP diagram:} IS shifts right in $t$, $t+1$ (economy moves up and right on the MP curve); returns to original in $t+2$ (economy moves back down, then below, the original intersection).

\noindent\textbf{Phillips Curve diagram:} The PC shifts up each period that $\widetilde{Y}>0$ (since $\pi_{t-1}$ increases). When $\widetilde{Y}<0$ the PC shifts back down. The economy traces a clockwise loop: up the initial PC, then left along shifted PCs, then back down.

\subsection*{(e) Numerical Computation [5 points]}

Given: $b_c = 0.6$, $b_T = 0.25$, $b_a = 0.1$, $b_i = 400$.

First: $b_c(1-b_T) = 0.6 \times 0.75 = 0.45$.

\textbf{Multiplier:}
\[
\text{With accelerator:}\quad \frac{1}{1-0.45-0.1} = \frac{1}{0.45} = \textcolor{sol}{\tfrac{20}{9} \approx 2.22}
\]
\[
\text{Without accelerator ($b_a=0$):}\quad \frac{1}{1-0.45} = \frac{1}{0.55} = \textcolor{sol}{\tfrac{20}{11} \approx 1.82}
\]

\textbf{Change in $Y$ for $\Delta R = 0.02$:}
\[
\text{With accelerator:}\quad \Delta Y = \frac{-b_i\,\Delta R}{1-b_c(1-b_T)-b_a} = \frac{-400 \times 0.02}{0.45} = \frac{-8}{0.45} = \textcolor{sol}{-\tfrac{160}{9} \approx -\$17.78\text{ billion}}
\]
\[
\text{Without accelerator:}\quad \Delta Y = \frac{-8}{0.55} = \textcolor{sol}{-\tfrac{160}{11} \approx -\$14.55\text{ billion}}
\]

\textbf{Conclusion:} The accelerator amplifies the interest rate's effect on output. A 2 percentage-point rise in $R$ reduces output by approximately \$3.23 billion more when $b_a = 0.1$ than when $b_a = 0$. This occurs because higher $R$ lowers $I$ directly, which lowers $Y$, which lowers $I$ further through the accelerator ($b_a \Delta Y$), which lowers $Y$ again. Each round of this feedback loop amplifies the total output contraction.

\subsection*{(f) Automatic Fiscal Stabilisers [6 points]}

\noindent\textbf{(i) Primary fiscal balance:}

Define $PB_t = b_T Y_t - \overline{G}$. At potential ($Y_t = \overline{Y}$), the budget is balanced: $\overline{G} = b_T\overline{Y}$, so:
\[
PB_t = b_T Y_t - b_T\overline{Y} = \textcolor{sol}{b_T(Y_t - \overline{Y})}
\]

Let $\hat{Y}_t \equiv Y_t - \overline{Y}$ denote the level output gap. Then $PB_t = b_T \hat{Y}_t$.

During a boom ($Y_t > \overline{Y}$, $\hat{Y}_t > 0$): $PB_t > 0$ — \textbf{the government runs a surplus}. Tax revenues rise automatically with income, while spending $\overline{G}$ is unchanged.

During a recession ($Y_t < \overline{Y}$, $\hat{Y}_t < 0$): $PB_t < 0$ — \textbf{the government runs a deficit}.

\noindent\textbf{(ii) Time path of $PB_t$:}

Since $PB_t = b_T(Y_t - \overline{Y})$, the fiscal balance mirrors the output gap scaled by $b_T$.

\begin{itemize}
  \item Periods $t$ and $t+1$ (boom, $\hat{Y} > 0$): $PB > 0$ — surplus, rising as the boom peaks.
  \item Period $t+2$ onward (correction, $\hat{Y} < 0$): $PB < 0$ — deficit, gradually recovering toward zero as $\hat{Y} \to 0$.
  \item Long run: $PB \to 0$ as $Y_t \to \overline{Y}$.
\end{itemize}

\noindent\textbf{Economic intuition:} Without any deliberate legislative change, the tax system automatically ``leans against'' the business cycle. During a boom, rising incomes push up tax revenues ($b_T Y_t$ rises) while government purchases remain at $\overline{G}$: the resulting surplus drains purchasing power and dampens the expansion. During a recession, falling incomes erode revenues while spending is unchanged: the deficit injects purchasing power and cushions the contraction. A higher marginal tax rate $b_T$ strengthens the automatic stabiliser (the $PB$ swings are larger relative to the output gap), smoothing cycles more effectively.


%----------------------------------------------------------------------
\section*{Question IV: The Phillips Curve and Credibility [15 points]}
%----------------------------------------------------------------------

\noindent Hybrid PC: $\pi_t = \lambda\overline{\pi} + (1-\lambda)\pi_{t-1} + \overline{v}\,\widetilde{Y}_t + \overline{o}_t$.

\noindent Rewrite as an inflation-gap equation: $(\pi_t - \overline{\pi}) = (1-\lambda)(\pi_{t-1}-\overline{\pi}) + \overline{v}\,\widetilde{Y}_t + \overline{o}_t$.

\subsection*{(a) One-Time Supply Shock [10 points]}

\noindent Starting at LR equilibrium: $\pi_{t-1} = \overline{\pi}$, $\widetilde{Y} = 0$. Shock at period $t$: $\overline{o}_t = \overline{o} > 0$; $\overline{o}_{t+s} = 0$ for $s \geq 1$. CB holds $\widetilde{Y} = 0$ for all periods.

\noindent\textbf{i. Closed-form expression for $(\pi_{t+s} - \overline{\pi})$:}

\textbf{Period $t$:}
\[
\pi_t - \overline{\pi} = (1-\lambda)\underbrace{(\pi_{t-1}-\overline{\pi})}_{=0} + \overline{v}\cdot 0 + \overline{o} = \overline{o}
\quad \text{(same for all } \lambda\text{)}
\]

\textbf{Period $t+s$} ($s \geq 1$, $\overline{o} = 0$, $\widetilde{Y} = 0$):
\[
\pi_{t+s} - \overline{\pi} = (1-\lambda)(\pi_{t+s-1} - \overline{\pi})
\]

Solving the recursion: $\quad \textcolor{sol}{\pi_{t+s} - \overline{\pi} = (1-\lambda)^s\,\overline{o}}$

As $s \to \infty$:
\begin{itemize}
  \item $\lambda = 0$: $(1-\lambda)^s = 1$ for all $s \Rightarrow$ inflation gap stays at $\overline{o}$ forever. Inflation is \textbf{permanently} stuck above target.
  \item $\lambda = 1$: $(1-\lambda)^s = 0$ for $s \geq 1 \Rightarrow$ inflation returns to $\overline{\pi}$ \textbf{immediately} in period $t+1$.
  \item $0 < \lambda < 1$: $(1-\lambda)^s \to 0$ exponentially. Inflation decays back to $\overline{\pi}$ at rate $(1-\lambda)$ per period.
\end{itemize}

\noindent\textbf{ii. Graph (time path of inflation):}

All three cases start at $\overline{\pi} + \overline{o}$ in period $t$.

\begin{itemize}
  \item $\lambda = 1$: drops back to $\overline{\pi}$ immediately in $t+1$ (one-period blip).
  \item $0 < \lambda < 1$: decays smoothly and monotonically toward $\overline{\pi}$ (exponential decay).
  \item $\lambda = 0$: remains flat at $\overline{\pi} + \overline{o}$ forever (horizontal line above target).
\end{itemize}

\textbf{Key qualitative difference:} When $\lambda > 0$, the shock is self-correcting --- inflation returns to target automatically even without any CB action or output gap. When $\lambda = 0$, the shock is permanently embedded in expectations and can only be undone through a deliberate disinflation (engineering $\widetilde{Y} < 0$), which is costly in terms of output.

\subsection*{(b) Credibility [5 points]}

In the model, $\lambda$ measures the fraction of agents whose inflation expectations are anchored to the announced target $\overline{\pi}$, rather than formed adaptively from past inflation. A high $\lambda$ is the quantitative expression of \textbf{credibility}: agents believe the Bank of Canada will return inflation to 2\% and therefore set wages and prices consistent with that belief.

\textbf{Why credibility matters:} When $\lambda$ is high, a supply shock such as higher oil prices or supply chain disruptions causes only a temporary spike in inflation. Workers do not demand large wage increases because they expect the price shock to dissipate; firms do not permanently raise prices for the same reason. The shock's inflationary effect self-corrects without requiring the CB to engineer a recession. When $\lambda$ is low, every shock becomes embedded in wage and price expectations, requiring painful tightening to reverse.

\textbf{Factors that lower $\lambda$:}
\begin{itemize}
  \item Past episodes of high inflation or missed targets (history undermines credibility).
  \item Political pressure on the central bank or perceived lack of independence.
  \item Changing or ambiguous inflation objectives.
  \item Central banks with a short institutional track record.
\end{itemize}

\textbf{How to raise $\lambda$:}
\begin{itemize}
  \item Formal inflation targeting mandates (explicitly committing to a numerical target, as the Bank of Canada does with its 1--3\% band).
  \item CB independence protected by legislation.
  \item Transparent communication (publishing forecasts, meeting minutes).
  \item A consistent track record of meeting targets over many years.
\end{itemize}


%----------------------------------------------------------------------
\section*{Question V: The AS-AD Model and Monetary Policy Trade-offs [15 points]}
%----------------------------------------------------------------------

\noindent Model: $\widetilde{Y} = \bar{a} - \tilde{b}(R-\overline{r})$ (IS), $R - \overline{r} = \overline{m}(\pi_t - \overline{\pi})$ (MP), $\pi_t = \pi^e_t + \overline{v}\,\widetilde{Y}_t + \overline{o}_t$ (PC), with $\pi^e_t = \pi_{t-1}$.

\subsection*{(a) The AD Curve [5 points]}

Substitute the MP rule into the IS curve:
\[
\widetilde{Y}_t = \bar{a} - \tilde{b}\,\overline{m}(\pi_t - \overline{\pi})
\]

\textcolor{sol}{
\[
\textbf{AD curve:} \quad \widetilde{Y}_t = \bar{a} - \tilde{b}\,\overline{m}\,(\pi_t - \overline{\pi})
\]
}

where $\tilde{b} = b_i/(1-b_c)$ and $\bar{a} = 0$ at long-run equilibrium (we allow $\bar{a} \neq 0$ when there are demand shocks). This is a downward-sloping relationship between $\widetilde{Y}$ and $\pi$.

\textbf{Why negative slope:} Higher inflation $\pi_t$ triggers the CB to raise $R$ (MP rule). Higher $R$ reduces investment and hence output $\widetilde{Y}$ (IS curve). Therefore, higher $\pi$ is associated with lower $\widetilde{Y}$: the AD curve slopes downward.

\subsection*{(b) Supply Shock with Three Policy Responses [10 points]}

\noindent Starting at LR equilibrium: $\widetilde{Y}=0$, $\pi = \overline{\pi}$.

\noindent AS at period $t$: $\pi_t = \pi_{t-1} + \overline{v}\,\widetilde{Y}_t + \overline{o} = \overline{\pi} + \overline{v}\,\widetilde{Y}_t + \overline{o}$ (shifted \textbf{up} by $\overline{o}$).

\noindent AD (baseline): $\widetilde{Y}_t = -\tilde{b}\,\overline{m}\,(\pi_t - \overline{\pi})$.

\medskip
\noindent\textbf{Response (i) --- Follow the rule:} AD unchanged.

Solving AS and AD simultaneously:
\begin{align*}
\pi_t - \overline{\pi} &= \overline{v}\,\widetilde{Y}_t + \overline{o} \\
\widetilde{Y}_t &= -\tilde{b}\,\overline{m}\,(\pi_t - \overline{\pi})
\end{align*}
Substituting: $(\pi_t - \overline{\pi})\bigl(1 + \overline{v}\tilde{b}\overline{m}\bigr) = \overline{o}$

\textcolor{sol}{
\[
\pi_t - \overline{\pi} = \frac{\overline{o}}{1 + \overline{v}\tilde{b}\overline{m}} > 0, \qquad \widetilde{Y}_t = \frac{-\tilde{b}\,\overline{m}\,\overline{o}}{1 + \overline{v}\tilde{b}\overline{m}} < 0
\]
}

Result (i): $\widetilde{Y}_t < 0$ and $\pi_t > \overline{\pi}$ --- \textbf{stagflation}: recession \emph{and} above-target inflation simultaneously.

\medskip
\noindent\textbf{Response (ii) --- Tighten to peg $\pi_t = \overline{\pi}$:}

From AS: $\overline{\pi} = \overline{\pi} + \overline{v}\,\widetilde{Y}_t + \overline{o} \Rightarrow \widetilde{Y}_t = -\overline{o}/\overline{v}$.

\textcolor{sol}{$\widetilde{Y}_t = -\overline{o}/\overline{v} < \widetilde{Y}_t^{(i)} < 0$, $\quad \pi_t = \overline{\pi}$.}

The CB achieves price stability but at the cost of a \textbf{deeper recession} than response~(i). The CB must shift AD leftward (raise $R$ aggressively beyond the rule) to force the economy up the new AS curve to the point $(\pi = \overline{\pi})$.

\medskip
\noindent\textbf{Response (iii) --- Ease to peg $\widetilde{Y}_t = 0$:}

From AS: $\pi_t = \overline{\pi} + \overline{o}$.

\textcolor{sol}{$\widetilde{Y}_t = 0$, $\quad \pi_t = \overline{\pi} + \overline{o} > \pi_t^{(i)} > \overline{\pi}$.}

The CB achieves output stability but accepts \textbf{higher inflation} than response~(i). The CB shifts AD rightward (cuts $R$ below the rule) to force the economy up the new AS curve to the point $(\widetilde{Y} = 0)$.

\medskip
\noindent\textbf{AS-AD diagram with all three outcomes:}

In the AS-AD space (horizontal axis $\widetilde{Y}$, vertical axis $\pi$): the AS curve shifts up by $\overline{o}$ (parallel shift). The three responses correspond to three different AD positions:
\begin{itemize}
  \item[(i)] Original AD (follows rule): equilibrium point $A_i = (\widetilde{Y}_i, \pi_i)$ with $\widetilde{Y}_i < 0$ and $\pi_i \in (\overline{\pi}, \overline{\pi}+\overline{o})$.
  \item[(ii)] AD shifts left (tighter): equilibrium point $A_{ii} = (-\overline{o}/\overline{v},\, \overline{\pi})$, further left on the new AS curve.
  \item[(iii)] AD shifts right (looser): equilibrium point $A_{iii} = (0,\, \overline{\pi}+\overline{o})$, at the top of the visible range on the new AS curve.
\end{itemize}

All three points lie on the \textbf{same shifted AS curve}. This is the key geometric insight: no AD position places the economy at $(\widetilde{Y}=0, \pi=\overline{\pi})$ because that point lies on the \emph{original} AS curve, not the shifted one. The CB faces a genuine trade-off.

\medskip
\noindent\textbf{Contrast with a demand shock:}

Suppose aggregate demand fell temporarily (IS shifted left). Under the CB's standard rule, the economy would settle at $(\widetilde{Y} < 0, \pi < \overline{\pi})$. By adjusting $R$ downward (shifting AD right), the CB can simultaneously restore $\widetilde{Y} = 0$ \emph{and} $\pi = \overline{\pi}$: the desired point $(\widetilde{Y}=0, \pi=\overline{\pi})$ lies on the AS curve (which has not shifted). A single policy action closes both gaps.

The supply shock destroys this ``free lunch.'' The shifted AS curve has moved the AS-AD equilibrium locus: there is no AD position consistent with both full employment and price stability. Any policy choice forces the CB to accept either a recession (to stabilise prices) or above-target inflation (to stabilise output). This is the fundamental challenge of \textbf{stagflation}.

The term for the simultaneous occurrence of a recession and above-target inflation is \textbf{stagflation} (stagnation + inflation).

\end{document}
